% preface.tex - Who this book is for and why
% ============================================================================

\chapter*{Preface}
\addcontentsline{toc}{chapter}{Preface}

\section*{Who This Book Is For}

This book is for scientists who need to program but never had formal training. You might be a chemistry grad student drowning in spectrometer data, a biologist tired of clicking through spreadsheets, or a physicist who inherited a folder of mysterious scripts from a graduated colleague. You know your science---you just need the programming skills to match.

This is not a computer science textbook. We won't discuss algorithm complexity or debate programming paradigms. Instead, we focus on what you need: reading data files, doing calculations, making figures, and automating the tedious parts of your work.

\section*{What You'll Learn}

By the end of this book, you'll be able to:

\begin{itemize}
    \item Set up a proper Python environment that won't break your computer
    \item Write scripts that read, process, and analyze your experimental data
    \item Create publication-quality figures
    \item Automate repetitive tasks that currently eat your afternoons
    \item Understand enough programming concepts to learn more on your own
\end{itemize}

We use chemistry examples throughout---titrations, spectra, kinetics---but the skills transfer to any scientific field. If you work with numbers and data, this book is for you.

\section*{How to Approach This Book}

Programming is a skill, like pipetting or playing piano. You can't learn it by reading---you have to do it. Every code example in this book is meant to be typed out and run. Copy-pasting works, but typing builds muscle memory, confidence, and forces you to notice details.

A reasonable pace is one chapter per week: read through, try all the examples, then tackle the exercises. Rushing won't help. Neither will reading ahead without practicing. When you hit something confusing, that's normal. Read it again, try it in Python, and it will click.

\section*{A Note on Errors}

You will see error messages. Many of them. This is not failure---it's how programming works. Every programmer, from beginner to expert, spends significant time reading error messages and figuring out what went wrong. The goal isn't to avoid errors; it's to get comfortable reading them and fixing the underlying problem.

When Python shows you an error, it's trying to help. The message tells you what went wrong and where. Learning to read these messages is as important as learning to write code.

\section*{Why Python?}

Python has become the default language for scientific computing, and for good reason. It's readable, has excellent libraries for data analysis and visualization, and has a massive community of scientists who share code and answer questions online. If you search for how to do almost anything in Python, you'll find answers.

Could you use R, MATLAB, or Julia? Sure. They're all fine languages. But Python's combination of general-purpose programming and scientific tools makes it the most practical choice for most researchers. Learn Python well, and you'll have a skill that serves you for your entire career.
